\begin{helper}[Strict PE growth from Algorithm 5.5 endpoint maximality]
Let $G$ be a finite loopless multigraph retaining individual edge
identities, let $[k]$ be the color set, and let $e_0$ be the distinguished
first edge.  Let $\mathcal E$ be a positive-rung ETT on a tree-sequence
$T$, with input coloring $\varphi_{\rm in}$, generating coloring
$\varphi$, and terminal extension kind $\Theta_n=PE$.  Fix a
partial coloring $\varphi_n$ and a tree-sequence $J$ carrying a
\noderef{CorrespondingExtendedTashkinovTree} certificate relative to
$\mathcal E$, with replacement ETT $\mathcal E_J$ generated under
$\varphi_n$.  Let $B_J$ be a \noderef{HierarchyBase} for $\mathcal E_J$
under $\varphi_n$, let $K_J$ be a \noderef{TashkinovHierarchy} on this
base with divider count zero, and let $\mathcal G_J$ be a
\noderef{GoodHierarchy} witness for $K_J$.  Assume
$J=B_J.\mathit{core}$.  Put
\[
 A=\noderef{HierarchyStarredStage}(K_J,0),\qquad
 D=\noderef{HierarchyStoredColorsAt}(K_J,0).
\]
Write $|W|$ for the number of displayed edges in a tree-sequence $W$, and
index those edges from zero.  For every tree-sequence $W$ for $G,e_0$
and every $\eta\in[k]$, put
\[
 C_W(\eta)=\left\{c\in[k]:\ \exists r\in\{0,\ldots,|W|-1\},\quad
 |A|\leq r<\min\bigl(
   \noderef{HierarchyFirstMissingPosition}(\varphi_n,W,\eta),|W|\bigr)
 \ \text{and}\ \varphi_n(W.\mathit{edge}(r))=c\right\}.
\]
This definition uses only the numerical lower cutoff $|A|$; $W$ need
not extend $A$ or $J$.  Let $m\in\mathbb N$ and let $U$ be any
tree-sequence for $G,e_0$ satisfying
\[
 |U|=|J|+m,
 \qquad m=0\Longrightarrow U=J.
\]
Assume the following exact one-step endpoint maximality condition:
\[
 \forall W\text{ a tree-sequence for }G,e_0,\quad
 \noderef{TashkinovAugmentation}(\varphi_n,U,W)
 \ \Longrightarrow
 \bigl(\forall\eta\in[k],\ \eta\in D\Longrightarrow
   \mathcal G_J.\Gamma^0_\eta\cap C_W(\eta)=\varnothing\bigr)
 \ \Longrightarrow\ \bot.
\]
Then $|J|<|U|$.
\end{helper}
\begin{proof}
Write
$D=\noderef{HierarchyStoredColorsAt}(K_J,0)$ and
$R=\bigcup_{\eta\in D}\mathcal G_J.\Gamma^0_\eta$.
The rung, series-presentation, and extension-kind agreements in
\noderef{CorrespondingExtendedTashkinovTree} transfer $\Theta_n=PE$
from $\mathcal E$ to $\mathcal E_J$.  The PE clause of
\noderef{GoodHierarchy} therefore says that
$J=B_J.\mathit{core}$ is not $R^-$-closed.  Expanding
\noderef{TreeColorSetMinusClosed}, choose a color
$c\in\noderef{TreeMissingColors}(\varphi_n,J)-R$ and an exact edge
$f\in\noderef{BoundaryColorEdges}(G,\varphi_n,V(J),c)$.
Unfolding \noderef{BoundaryColorEdges} and \noderef{BoundaryEdges} gives
vertices $x\in V(J)$ and $y\notin V(J)$ with
$G.\mathit{ends}(f)=\{x,y\}$.

Construct an ordered tree $W=J+f$ using the data of
\noderef{TreeSequence}, with positive length $|J|+1$: copy every old
vertex and edge entry, append $y$ as the last vertex, and append the
individual identity $f$ as the last
edge.  The old entries retain their injectivity.  The new vertex is fresh
because $y\notin V(J)$.  The new edge is also fresh: every old displayed
edge has both endpoints among the old vertices by the attachment field,
whereas $f$ has the endpoint $y$ outside that set.  The old attachment
equations and distinguished first edge copy, and the boundary endpoint
equation supplies the final attachment to the occurrence of $x$ in $J$.
Thus $W$ is an actual exact ordered tree-sequence, $J$ is its exact
\noderef{TreeSequenceInitialSegment}, and
$|W|=|J|+1$.  Since $\varphi_n(f)=c$ and $c$ is missing on $J$, these are
exactly the fields of a
\noderef{TashkinovAugmentation} from $J$ to $W$.

Because the divider count is zero,
\noderef{HierarchyStarredStage}$(K_J,0)=B_J.\mathit{core}=J$.  Hence the
only possible displayed edge contributing to $C_W(\eta)$ is the newly
appended edge at index $|J|$, and its color is $c$.  Whatever the
totalized \noderef{HierarchyFirstMissingPosition} of $\eta$ in $W$,
we therefore have $C_W(\eta)\subseteq\{c\}$.  The choice $c\notin R$
says that $c$ belongs to no Gamma pair indexed by $D$.  Therefore
$\mathcal G_J.\Gamma^0_\eta\cap C_W(\eta)=\varnothing$ for every
$\eta\in D$, so $W$ satisfies every avoidance condition in the
endpoint maximality hypothesis.

Suppose $m=0$.  The endpoint hypothesis gives $U=J$, so the preceding
augmentation becomes an admissible one-step augmentation of $U$.  This
contradicts endpoint maximality.  Thus $m>0$, and the length formula gives
$|U|=|J|+m>|J|$, as required.
\end{proof}
